\section{Preliminaries}


\subsection{Notation}
\label{sec:notation}
Let \(\Field\) be a small finite field of order \(q\), and let
\(\ExtensionField=\mathsf{GF}(q^\kappa)\) be a degree-\(\kappa\)
extension of \(\Field\). Thus \(q=|\Field|\) is a prime power
and \(\Field\subset\ExtensionField\).
For security parameter \(\lambda\), we choose \(\kappa\) to
satisfy the soundness bounds of our protocols, including
\(
    \lvert\ExtensionField\rvert
    =
    q^\kappa
    \geq
    2^\lambda.
\)
We write \(N\) for the witness length over \(\Field\) and set
\(m=N/\kappa\), assuming \(\kappa\mid N\).
For the main circuit instantiations, we require \(q>\kappa\).
This condition concerns the field order \(q\).
The intermediate-field Hadamard variant in
Section~\ref{sec:hadamard} also treats \(q<\kappa\); its tabulated
implementation reports online time, table construction, and total
storage separately.
Our main instantiations use odd prime fields;
Section~\ref{sec:binary-packing} treats binary fields using a small
extension as the working field.

We use lowercase letters for scalars, bold lowercase letters for
column vectors, and bold uppercase letters for matrices.
An overline denotes an element, vector, or matrix over
\(\ExtensionField\), as in \(\ext a\), \(\ext{\bm a}\), and
\(\ext{\bm A}\). Transposition is denoted by \(\top\).

Let
\(\overline{\beta}_1,\ldots,\overline{\beta}_\kappa\)
be a basis of \(\ExtensionField\) over \(\Field\), and let
\(\overline{\beta}_1^\dagger,\ldots,
\overline{\beta}_\kappa^\dagger\)
denote its trace-dual basis.
The polynomial rings over \(\Field\) and \(\ExtensionField\) are
denoted by \(\Field[X]\) and \(\ExtensionField[X]\), respectively.
We use \(X,Y,Z\) for formal indeterminates,
\(\overline{x},\overline{y},\overline{z}\) for extension-field
challenges, and \(i,j,\ell\) for indices.
% Sets are denoted using blackboard bold font,\eg,
% The set of natural numbers is $\mathbb{N}$,
% the set of powers of two is $2^\mathbb{N}$,
% and the set of real numbers is $\mathbb{R}$.
We write $[n]=\{1,\ldots,n\}$ and use $|\cdot|$ for cardinality.


We also use the following matrix and vector operations:
\begin{itemize}[noitemsep, topsep = 0pt]
\item inner product $c = \langle \bm a, \bm b \rangle$; addition ($\bm C = \bm A + \bm B$); transposition ($\bm C = \bm A^\top$);
\item multiplication
($\bm C = \bm A\bm B$); Kronecker (tensor) product ($\bm C = \bm A \otimes \bm B$);
\item Hadamard (element-wise) product $\bm{A} \circ \bm{B}$;
\item concatenation of vectors $\bm{c} = (\bm c_1 \| \bm c_2 \|\cdots \| \bm c_\kappa)$; and
\item horizontal packing of column vectors $\bm C = [\bm c_1 \mid \bm c_2 \mid \cdots \mid \bm c_\kappa]$.
\end{itemize}
    
We charge \(\bb\), \(\be\), and \(\ee\) for a base-field,
mixed-field, and extension-field multiplication, respectively,
and absorb additions into these costs. Schoolbook extension-field
arithmetic gives
\(
    \be=\kappa\bb
\) and
\(
    \ee=\kappa\be=\kappa^2\bb
\).
Tabulated implementations additionally allow random-access tables,
charging one unit per base-field coordinate read or written.
We distinguish online arithmetic, public table construction,
and total storage, including read-only tables.


We denote the prover and verifier by $\Prover$ and $\Verifier$.
We specify relations as
$\Relation=\{\text{public input}:\text{private input}\mid\text{constraint}\}$.

\subsection{SNARK Protocol in Vector Form}
\label{sec:defgeneral}
We express circuit satisfiability as vector identities and represent
multilinear-extension (MLE) evaluations as inner products with
tensor vectors. This formulation exposes the matrix structure used
by our packing reductions. We first recall the PLONK constraint
system and then state the evaluation and proof interfaces.

\begin{definition}[PLONK constraint system for circuit
satisfiability~\cite{iacr/GabizonWC19,iacr/AmbronaSTW22}]
\label{def:plonk}
An arithmetic circuit with \(N\) gates and \(n_{\mathsf{io}}\)
public inputs and outputs is specified by five selector vectors
\(
    \bm q_L,\bm q_R,\bm q_M,\bm q_O,\bm q_C
    \in\Field^N
\)
and an augmented wiring permutation
\(
    \sigma:
    [3N+n_{\mathsf{io}}]
    \to
    [3N+n_{\mathsf{io}}],
\)
encoding wire-copy constraints and connections to the public inputs
and outputs.
Given a public-I/O vector
\(\bm{\mathsf{io}}\in\Field^{n_{\mathsf{io}}}\), let
\(\bm w_L,\bm w_R,\bm w_O\in\Field^N\) denote the left, right,
and output wire assignments, respectively, and define the
augmented assignment
\(
    \bm w
    \coloneqq
    (\bm w_L\mathbin{\|}\bm w_R\mathbin{\|}\bm w_O
    \mathbin{\|}\bm{\mathsf{io}})
    \in\Field^{3N+n_{\mathsf{io}}}.
\)
The wire assignment satisfies the circuit when:
\begin{align}
    \bm q_L\circ\bm w_L
    +\bm q_R\circ\bm w_R
    +\bm q_M\circ(\bm w_L\circ\bm w_R)
    +\bm q_O\circ\bm w_O
    +\bm q_C
    &=\bm 0_N,
    \label{eq:plonk1}\\
    \sigma(\bm w)
    &=\bm w.
    \label{eq:plonk2}
\end{align}
Here,
\(\sigma(\bm w)\) denotes the coordinate permutation of
\(\bm w\) induced by \(\sigma\).
\end{definition}

\begin{definition}[R1CS constraint system for circuit satisfiability]
\label{def:r1cs}
An arithmetic circuit is specified by three sparse matrices
\(
    \bm S_L,\ \bm S_R,\ \bm S_O \in \Field^{N\times N},
\)
where \(N\) is the total number of multiplication gates and I/O gates
(the public inputs and outputs). Let
\(
    \bm w = (\bm w^{\mathsf{mul}} \,\|\, \bm{\mathsf{io}}) \in \Field^{N}
\)
be the assignment to the multiplication gates and the I/O gates.
The circuit is satisfied if and only if:
\[
    (\bm S_L \bm w) \circ (\bm S_R \bm w) = \bm S_O \bm w .
\]
\end{definition}

The following tensor representation identifies an MLE query with a
linear functional on its evaluation vector.

\begin{definition}[Tensor span]
\label{def:tensorspan}
Let \(n \in 2^\mathbb{N}\) and let
\(
    \bm x
    =
    (x_1,\ldots,x_{\log_2 n})
    \in\Field^{\log_2 n}.
\)
The \emph{tensor span} of \(\bm x\) is the length-\(n\) vector
\(
    \tensorspan{\bm x}
    \coloneqq
    \bigotimes_{i=1}^{\log_2 n} (1-x_i, x_i)^\top
    \in\Field^n.
\)
The same construction applies to
\(\ext{\bm x}\in\ExtensionField^{\log_2 n}\).
We define the sets of tensor spans over \(\Field\) and
\(\ExtensionField\), respectively, by:
\begin{align*}
    \TensorSpace(n)
    &\coloneqq
    \left\{
        \tensorspan{\bm x}
        \;\middle|\;
        \bm x\in\Field^{\log_2 n}
    \right\}
    \subseteq\Field^n,
    \qquad
    \ext{\TensorSpace}(n)
    \coloneqq
    \left\{
        \tensorspan{\ext{\bm x}}
        \;\middle|\;
        \ext{\bm x}\in\ExtensionField^{\log_2 n}
    \right\}
    \subseteq\ExtensionField^n.
\end{align*}
\end{definition}

\paragraph{Products of tensor weights.}
We also allow tensor queries with arbitrary two-entry factors.
Coordinatewise products preserve this representation:
\begin{equation}
    \tensorspan{\ext{\bm z}}\circ\tensorspan{\ext{\bm y}}
    =\bigotimes_{j=1}^{\log_2n}
    \begin{pmatrix}
        (1-\ext z_j)(1-\ext y_j)\\
        \ext z_j\ext y_j
    \end{pmatrix}.
    \label{eq:tensor-weight-product}
\end{equation}
When \((1-\ext z_j)(1-\ext y_j)\neq0\) for every \(j\),
factoring out these entries gives:
\[
    \tensorspan{\ext{\bm z}}\circ\tensorspan{\ext{\bm y}}
    =\left(\prod_j(1-\ext z_j)(1-\ext y_j)\right)
      \bigotimes_j
      \begin{pmatrix}
          1\\
          \dfrac{\ext z_j}{1-\ext z_j}
          \dfrac{\ext y_j}{1-\ext y_j}
      \end{pmatrix}.
\]
To express the result as a scaled MLE query, set
\(\ext d_j=(1-\ext z_j)(1-\ext y_j)+\ext z_j\ext y_j\).
If every \(\ext d_j\neq0\), then:
\[
    \tensorspan{\ext{\bm z}}\circ\tensorspan{\ext{\bm y}}
    =\left(\prod_j\ext d_j\right)
      \tensorspan{(\ext z_j\ext y_j/\ext d_j)_j}.
\]
Our protocols retain the two-entry factors in
Eq.~\eqref{eq:tensor-weight-product}, which also covers zero
entries and vanishing normalization factors without division.

\begin{lemma}[Schwartz-Zippel-DeMillo-Lipton Lemma]
\label{lemma:schwartz}
Let
\( \phi(\bm{x}) : \Field^\upsilon \to \Field \) be a nonzero \( \upsilon \)-variate polynomial such that the degree of \( \phi \) in each variable is at most \( d \). Let \( \bm{x} \ChooseRandom \mathbb{X}^{\upsilon} \) for some finite set \( \mathbb{X} \subseteq \Field \). Then:
\[
    \Pr_{\bm{x} \ChooseRandom \mathbb{X}^{\upsilon}}\!\left[ \phi(\bm{x}) = 0 \right]
    \le \frac{\upsilon \cdot d}{|\mathbb{X}|}.
\]
\end{lemma}

\begin{corollary}[Zero check in matrix form]
\label{corollary:schwartz}
Let \(n\) be a positive power of two and let
\(
    \bm a\in\Field^n\setminus\{\bm 0_n\}.
\)
Viewing \(\bm a\) as a vector over \(\ExtensionField\), sample
\(
    \ext{\bm x}
    \ChooseRandom
    \ExtensionField^{\log_2 n}.
\)
Then:
\[
    \Pr_{\ext{\bm x}\ChooseRandom
        \ExtensionField^{\log_2 n}}
    \!\left[
        \left\langle
            \bm a,
            \tensorspan{\ext{\bm x}}
        \right\rangle
        =0
    \right]
    \leq
    \frac{\log_2 n}{\lvert\ExtensionField\rvert}.
\]
\end{corollary}

\begin{definition}[Inner-product proof]
\label{def:ip}
Let \(m\) be a positive power of two. Suppose that the prover
\(\Prover\) is bound to two vectors
\(\ext{\bm a},\ext{\bm b}\in\ExtensionField^m\), and that the
verifier \(\Verifier\) holds the corresponding commitments.
Given a claimed value \(\ext{v}\in\ExtensionField\), an
\emph{inner-product proof} is an interactive protocol for the
relation:
\[
    \Language_{\mathsf{ip}}
    \coloneqq
    \left\{
        \ext{v}\in \ExtensionField,~      \ext{\bm a},\ext{\bm b}\in\ExtensionField^m
        \;\middle|\;
        \begin{array}{l}
            \ext{v}
            =\langle\ext{\bm a},\ext{\bm b}\rangle
        \end{array}
    \right\}.
\]
\end{definition}

Sum-check~\cite{crypto/Thaler13} and folding~\cite{asiacrypt/CCYY25}
give the following reduction to two MLE claims. We represent
sum-check round polynomials by their coefficients; the round check
\(g(0)+g(1)=c\) then applies in every characteristic.

\begin{fact}[\bf Inner-product reduction]
\label{fact:ip}
There is an inner-product reduction with prover cost
\(4m+O(\log m)\) multiplications in \(\ExtensionField\), communication
\(2\log_2 m+2\) $\ExtensionField$-elements, and verifier cost
\(O(\log m)\) $\ExtensionField$-multiplications.
Over \(\log_2 m\) rounds, the verifier samples challenges collected
in
\(
    \ext{\bm x}\in\ExtensionField^{\log_2 m},
\)
and the protocol reduces the inner-product relation to two
terminal MLE claims:
\[
    \ext{v}_1
    \EqOrNot
    \left\langle
        \ext{\bm a},
        \tensorspan{\ext{\bm x}}
    \right\rangle
    \quad\wedge\quad
    \ext{v}_2
    \EqOrNot
    \left\langle
        \ext{\bm b},
        \tensorspan{\ext{\bm x}}
    \right\rangle,
\]
where the prover supplies \(\ext{v}_1,\ext{v}_2\in\ExtensionField\).
The reduction has soundness error at most
\(
    \frac{2\log_2 m}{\lvert\ExtensionField\rvert}.
\)
\end{fact}

\begin{definition}[Multilinear polynomial commitment scheme in matrix form]
\label{def:pcs}
Let \(n \in 2^\mathbb{N}\). A multilinear PCS over \(\Field\)
with evaluations over \(\ExtensionField\) consists of algorithms
$(\Setup,\Com)$ and an opening protocol $\Open$:
\begin{itemize}[leftmargin=*,labelsep=0.5em]
    \item
    \(
        \pp\leftarrow\Setup(1^\lambda,n)
    \)
    takes the security parameter and the committed vector length
    \(n\), and outputs public parameters \(\pp\).

    \item
    \(
        C\leftarrow\Com(\pp,\bm a;\rho)
    \)
    takes a vector \(\bm a\in\Field^n\) and fresh randomness
    \(\rho\), and outputs a commitment \(C\). We write
    \(\Com(\bm a;\rho)\) when \(\pp\) is clear from context.

    \item The opening protocol
    \(
        \mathsf{bit}
        \leftarrow
        \Open\!\left(
            \pp,C,\ext{v},\ext{\bm x}
            :
            \bm a,\rho
        \right)
    \)
    is an interactive argument between \(\Prover\) and
    \(\Verifier\) for the relation:
    \[
        \Relation_{\mathsf{open}}
        \coloneqq
        \left\{
            \bigl(
                \pp,C,\ext{v},\ext{\bm x}
            \bigr)
                :
            \bigl(
                \bm a,\rho
            \bigr)
            \;\middle|\;
            \begin{aligned}
                C=\Com(\pp,\bm a;\rho) \wedge
                \ext{v}=\left\langle
            \bm a,
            \tensorspan{\ext{\bm x}}
        \right\rangle
            \end{aligned}
        \right\}.
    \]
    At the end of the protocol, \(\Verifier\) outputs
    \(\mathsf{bit}\in\{\TRUE,\FALSE\}\).
\end{itemize}
\end{definition}



\subsection{Extension-Field Arithmetic}

\begin{fact}[Base-Field Embedding]
\label{fact:injective}
Let\/ \(\ExtensionField/\Field\) be an extension.
The canonical inclusion
\(
    \iota:\Field\hookrightarrow\ExtensionField
\)
is injective and induces an injective inclusion
\(
    \Field[X]\hookrightarrow\ExtensionField[X].
\)
\end{fact}

The field trace maps extension-field elements back to \(\Field\)
and gives an explicit formula for their base-field coordinates.
\begin{definition}[Extension-field trace]
\label{def:trace}
The \emph{field trace} from \(\ExtensionField = \mathsf{GF}(q^\kappa)\) to \(\Field = \mathsf{GF}(q)\)
is the \(\Field\)-linear map:
\[
    \trace:\ExtensionField\to\Field,
    \qquad
    \ext{\alpha}
    \longmapsto
    \trace(\ext{\alpha})
    \coloneqq \ext{\alpha} + \ext{\alpha}^q + \ext{\alpha}^{q^2} + \cdots + \ext{\alpha}^{q^{\kappa-1}}.
\]
In particular, for all \(a,b\in\Field\) and
\(\ext{\alpha},\ext{\beta}\in\ExtensionField\):
\[
    \trace\!\left(
        a\ext{\alpha}+b\ext{\beta}
    \right)
    =
    a\trace(\ext{\alpha})
    +
    b\trace(\ext{\beta}).
\]
\end{definition}

\begin{definition}[Basis and trace-dual basis]
\label{def:basis}
An ordered tuple
\(
    \mathcal B
    =
    \bigl(
        \ext{\beta}_1,\ldots,\ext{\beta}_\kappa
    \bigr)
    \in\ExtensionField^\kappa
\)
is a \emph{basis} of \(\ExtensionField\) over \(\Field\) if every
\(\ext{\alpha}\in\ExtensionField\) admits a unique representation:
\[
    \ext{\alpha}
    =
    \sum_{i=1}^{\kappa}a_i\ext{\beta}_i,
    \qquad
    a_1,\ldots,a_\kappa\in\Field.
\]
The \emph{trace-dual basis} of \(\mathcal B\) is the unique basis
\(
    \mathcal B^\dagger
    =
    \bigl(
        \ext{\beta}_1^\dagger,\ldots,
        \ext{\beta}_\kappa^\dagger
    \bigr)
\)
satisfying:
\[
    \trace\!\left(
        \ext{\beta}_i\ext{\beta}_j^\dagger
    \right)
    =
    \delta_{ij}
    \qquad
    \text{for all }i,j\in[\kappa],
\]
where \(\delta_{ij}\) is the Kronecker delta. Consequently, the
coordinates of \(\ext{\alpha}\) with respect to \(\mathcal B\)
can be recovered as:
\[
    a_i
    =
    \trace\!\left(
        \ext{\alpha}\ext{\beta}_i^\dagger
    \right),
    \qquad
    i\in[\kappa].
\]
\end{definition}

Our base-field instantiations use interpolation nodes in \(\Field\), so that the Lagrange
basis has base-field coefficients even when the interpolated
polynomial has extension-field values. The same definition over
an intermediate field \(\mathbb K\) is used for the small-field
variant in Section~\ref{sec:hadamard}.


\begin{definition}[Lagrange basis over the extension field]
\label{def:lagrange}
Assume that \(q\geq\kappa\), and fix an ordered set of distinct
nodes \(\mathcal H=\{\alpha_1,\ldots,\alpha_\kappa\}\subseteq\Field\).
The indices \(j\in[\kappa]\) label the chosen field elements.
The vanishing polynomial of \(\mathcal H\) is:
\[
    \Psi(X)
    \coloneqq
    \prod_{j=1}^{\kappa}(X-\alpha_j)
    \in\Field[X].
\]
For each \(j\in[\kappa]\), define the Lagrange basis polynomial
\(
    \psi_j(X)
    \coloneqq
    \prod_{\substack{\ell\in[\kappa]\\ \ell\neq j}}
    \frac{X-\alpha_\ell}{\alpha_j-\alpha_\ell}
    \in\Field[X].
\)
These polynomials satisfy:
\[
    \psi_j(\alpha_i)=\delta_{ij}
    \qquad
    \text{for all }i,j\in[\kappa].
\]
Moreover, \(\{\psi_j(X)\}_{j=1}^{\kappa}\) forms an
\(\ExtensionField\)-basis for the space of polynomials in
\(\ExtensionField[X]\) of degree less than \(\kappa\).
Consequently, every \(g(X)\in\ExtensionField[X]\) with
\(\deg g<\kappa\) admits the unique representation:
\[
    g(X)
    =
    \sum_{j=1}^{\kappa}
    \ext{a}_j\psi_j(X),
    \qquad
    \ext{a}_j=g(\alpha_j)\in\ExtensionField.
\]
\end{definition}
